\section{应用推广情况和经济社会效益}

\subsection{成果形态与应用范围}

本项目围绕城市复杂场景感知形成了软件、平台和系统等多种成果形态，将动态视觉极限感知、多源协同融合和复杂任务智能理解等技术转化为可部署、可集成的行业解决方案。相关成果已面向智慧交通、智能制造、智能安防和智慧应急等领域开展规模化应用，覆盖交通基础设施运行监测、电子制造智能质检、车辆充电安全监管、城市消防安全感知等典型业务。通过与行业用户既有设备、数据平台和业务流程相结合，项目成果实现了从关键算法、功能模块到综合平台的逐级落地。

项目技术的应用对象既包括城市道路、桥梁、隧道、边坡和交通枢纽等大范围基础设施，也包括工业生产线上的锡膏、印制电路板和电子屏幕等精细检测对象，还包括电动车、充电桩、消防设施等数量庞大的物联网终端。上述场景在目标尺度、采样频率、运行环境和实时性要求上差异显著，验证了项目技术面向多行业、多设备和多任务进行适配的能力，并为后续跨区域复制推广奠定了应用基础。

从部署形态看，项目成果并非以单一算法原型进入应用，而是根据用户业务基础形成不同层级的交付方式。对于已经具备传感器、生产设备或物联网终端的用户，项目以感知算法、智能质检模块和风险识别模块等软件能力嵌入既有系统；对于需要统一管理大量设施和节点的用户，项目进一步建设数据接入、状态展示、异常预警与业务联动平台；对于交通基础设施监测和智慧消防等跨区域任务，则将前端设备、通信接入、融合分析与指挥处置组织为完整系统。软件、平台和系统三类成果相互衔接，使同一技术体系能够适配不同建设基础和投资规模。

从技术供给与行业需求的对应关系看，动态视觉和精细目标感知能力主要解决高速目标、微小缺陷与弱小异常能否被及时发现的问题，多源协同融合能力主要解决跨设备、跨区域信息能否统一接入和可信汇聚的问题，多层次认知与推理能力则服务于风险研判、任务规划和协同处置。三类能力在行业应用中不再作为彼此独立的研究模块出现，而是围绕用户的质量、安全、效率和管理目标进行组合，由此形成面向交通、制造、安防和应急的差异化解决方案。

\subsection{成果转化与产业化路径}

项目成果遵循“关键技术研发---产品功能封装---行业平台集成---典型场景验证---规模化推广”的转化路径。首先，将极限感知、协同融合和认知推理方法封装为可调用的软件模块和算法组件；其次，根据生产线、道路设施、物联网终端和消防设备等对象建立数据接口与业务接口；再次，在行业平台中完成状态监测、异常识别、风险预警和任务处置等功能集成；最后，通过实际用户的长期运行验证技术效果，并将成熟模块复制到同类设备、生产线或区域平台。

这一转化路径兼顾了技术先进性与工程可实施性。一方面，核心算法可以随研究进展独立升级，使产品持续获得新的感知和理解能力；另一方面，平台保留与既有设施、终端和业务系统的连接关系，避免每次算法更新都重新建设完整系统。在智能制造场景中，算法能力可以嵌入不同检测工位；在智慧交通、智能安防和智慧应急场景中，平台可以通过扩展接入节点和业务模块扩大覆盖范围，从而降低成果由单点示范走向规模化应用的组织成本。

项目现有应用已经覆盖企业生产、城市基础设施运行和公共安全治理等不同类型业务，说明成果转化并不依赖单一行业的市场需求。制造业应用通过提高检测能力和产品良率形成较为直接的经营收益；交通、安防和应急应用则通过平台建设、运行服务、风险预警和损失避免形成综合效益。多种效益实现方式共同构成项目产业化基础，也增强了成果面对行业周期和场景变化时的持续推广能力。

\subsection{总体经济效益}

项目成果已在多家完成单位实现产业化应用。按照项目应用单位汇总口径，累计形成销售收入1379.52亿元、直接经济效益72.29亿元。其中，创维、天珑、深城交和深智城分别依托智能制造、城市交通和智慧城市等相关业务推动成果转化，形成了从技术研发、产品集成到行业应用的产业化链条。具体情况见表\ref{tab:economic-benefits}。

\begin{table}[htbp]
  \centering
  \caption{四家完成单位累计经济效益（金额单位：亿元）}
  \label{tab:economic-benefits}
  \renewcommand{\arraystretch}{1.25}
  \begin{tabular}{|c|c|r|r|}
    \hline
    \textbf{应用单位} & \textbf{统计周期} & \textbf{销售收入} & \textbf{直接经济效益} \\
    \hline
    创维   & 2023--2025年 & 684.74  & 8.90  \\
    \hline
    天珑   & 2022--2025年 & 633.04  & 58.02 \\
    \hline
    深城交 & 2023--2025年 & 39.40   & 1.85  \\
    \hline
    深智城 & 2023--2025年 & 22.34   & 3.52  \\
    \hline
    累计   & ---           & 1379.52 & 72.29 \\
    \hline
  \end{tabular}
\end{table}

表中销售收入体现应用单位相关业务的市场规模，直接经济效益体现项目技术应用所形成的新增价值。两类指标的统计含义不同，不作简单相加。表中分项数据按展示精度保留两位小数，累计值采用项目材料的统一汇总口径。

创维的应用证明显示，在2023--2025年统计周期内，企业累计建设集自动化、数字化和人工智能质检于一体的高科技生产线27条，产线良品率提升1.3\%，形成直接经济效益约8.90亿元。相关成果说明，项目技术不仅完成了实验室条件下的缺陷识别验证，而且进入了连续运行的生产过程，并通过减少缺陷漏检、稳定质量控制和支撑产线自动化形成可量化收益。对应汇总口径下，创维形成销售收入684.74亿元、直接经济效益8.90亿元。

天珑的应用证明显示，在2022--2025年统计周期内，企业依托人工智能算法与生物启发式智能感知技术新建自动化产线16条，相关产线对应的销售增量占比达到80\%，直接经济效益约58.02亿元。与创维的27条产线合计，项目相关智能制造技术共应用于43条生产线，形成了从智能终端制造到智能家电制造的多类型产线应用。对应汇总口径下，天珑形成销售收入633.04亿元、直接经济效益58.02亿元。

深城交主要承担项目成果在城市交通基础设施监测和交通运行服务中的转化应用，在2023--2025年统计周期内形成销售收入39.40亿元、直接经济效益1.85亿元。与制造企业以生产线质检和产品经营收益为主的效益模式不同，交通领域的价值更多通过平台建设、设施运行监测、数据服务和安全管养支撑实现。

深智城主要推动项目成果在智慧城市平台和城市运行服务中的转化应用，在2023--2025年统计周期内形成销售收入22.34亿元、直接经济效益3.52亿元。四家单位分别代表智能制造产品化、城市交通平台化和智慧城市服务化等主要转化路径，共同构成项目72.29亿元直接经济效益的应用基础。

从效益构成看，智能制造是直接经济效益的主要来源之一。相关技术进入全球领先的智能终端全栈服务商和国内智能家电领军企业的生产环节，应用于43条生产线，覆盖锡膏三维智能检测、PCB装配缺陷智能检测和电子屏幕划痕检测等任务，在提升产品良率的同时形成直接经济效益超过66亿元。智慧交通、智能安防和智慧应急等应用则通过平台建设、设施监测、风险预警与运维服务形成经济效益，并进一步体现为基础设施安全运行、事故风险降低和城市治理效率提升。

\subsection{主要用户与推广范围}

项目成果面向不同层级的政府部门、行业龙头企业和基层业务单位形成了较为广泛的用户基础。在智慧交通领域，主要用户包括广东、福建、甘肃、四川等省级交通运输主管部门及相关市级交通运输部门；在智能制造领域，成果服务于智能终端全栈服务商和国内智能家电领军企业；在智能安防领域，技术已应用于港口、机场等海关单位以及深圳市内10余个街道办事处；在智慧应急领域，成果已在广东省内10余个市级消防救援支队及基层消防单位推广，并面向卫生健康等相关部门形成应用。

上述用户覆盖省、市、街道等不同管理层级，也涵盖交通、制造、海关、消防和医疗卫生等不同业务体系。一方面，跨区域应用说明项目技术能够适应不同地区的设施条件与管理需求；另一方面，政府部门、行业企业和基层单位的共同参与，使项目成果能够贯通设施接入、状态监测、风险研判和业务处置等环节，形成面向实际业务闭环的推广模式。

从地域分布看，智慧交通成果已由深圳本地应用拓展至广东、福建、甘肃、四川等省份，智慧应急成果覆盖广东省内多个市级消防救援支队和基层消防单位，智能安防成果延伸至深圳市内10余个街道办事处。不同地区在基础设施规模、存量设备类型和业务管理流程上存在差异，跨区域部署要求项目平台具备多类型设备接入、异构数据融合和功能模块扩展能力。现有应用范围表明，项目成果已经由单个示范点逐步形成面向区域和行业的推广网络。

从用户性质看，项目同时服务政府管理部门、公共安全单位和市场化企业。政府及公共安全用户重点关注设施覆盖、运行可靠、风险预警和协同处置，制造企业重点关注检测精度、生产节拍、产品良率和经营收益。项目以统一的智能感知技术为基础，根据不同用户的评价目标调整系统形态和功能组合，使科研成果既能够支撑公共治理，也能够形成产业收益。这种多主体应用结构有利于形成需求反馈、技术迭代和规模推广之间的正向循环。

\subsection{典型应用成效}

\subsubsection{智慧交通基础设施运行监测}

在智慧交通领域，项目建成了接入设施体量全国最大的城市级交通基础设施运行监测平台，形成了覆盖多种典型场景和设施类型的应用示范。平台面向道路、桥梁、隧道、边坡和交通枢纽等设施开展持续监测，已实现道路8455公里、桥梁4180座、隧道170条、边坡3590座、交通枢纽8座的规模化监测应用。相关成果将分散的设施感知数据汇聚到统一平台，为设施状态识别、风险分析、安全管养和运行决策提供支撑，推动交通基础设施管理由人工巡检和分散处置向在线监测与综合研判转变。

交通基础设施具有数量多、分布广、类型复杂和运行连续等特点，单靠人工巡查难以持续获得全部设施状态。项目平台将道路、桥梁、隧道、边坡及枢纽等不同对象纳入统一监测体系，使管理人员能够从城市或区域尺度掌握设施运行情况，并针对异常状态开展进一步核查和处置。平台的价值不仅体现在接入数量上，还体现在对不同设施数据的统一组织，以及监测、分析和管养业务之间的衔接。

相关应用获得交通运输管理部门出具的应用证明，表明平台已用于实际交通基础设施养护与运行管理。大规模设施接入为项目带来平台建设和技术服务收益，同时能够减少重复建设和分散管理带来的资源浪费。随着接入设施类型增加，既有平台的数据接口、分析能力和管理流程可以继续复用，为其他城市或区域开展同类基础设施数字化监测提供可参照的实施路径。

\subsubsection{智能制造质量检测}

在智能制造领域，项目技术应用于智能终端和智能家电制造过程，形成锡膏三维智能检测、PCB装配缺陷智能检测和电子屏幕划痕检测等解决方案。相关技术已应用于43条生产线，通过对细微结构、装配缺陷和表面弱小异常进行自动识别，提升生产环节的检测一致性与缺陷发现能力，产品良率得到显著提升，并实现直接经济效益超过66亿元。该应用表明，项目在微小目标、精细结构和复杂背景条件下形成的感知能力，能够直接服务于制造过程质量控制和生产效率提升。

三类检测任务对应不同的质量控制难点。锡膏三维检测需要从细微几何形态中判断印刷质量，PCB装配检测需要在元器件密集、结构复杂的图像中发现装配异常，电子屏幕划痕检测则需要识别低对比度、弱纹理的表面缺陷。项目将极限场景感知与多源信息处理方法转化为生产线可用的检测功能，使原本依赖人工经验的细粒度判断能够进入自动化质量控制流程。

从产业化结果看，创维建设27条相关高科技生产线，天珑新建16条自动化产线，两类应用共同构成43条生产线的规模化部署。创维产线良品率提升1.3\%，说明智能检测不仅增加了缺陷发现手段，也对生产过程优化和最终产品质量产生了实际作用；天珑相关产线形成显著销售增量，说明技术应用能够与企业产能建设和市场经营形成联动。制造业场景由此成为项目科研成果转化为直接经济效益的核心载体。

\subsubsection{城市智能安防}

在智能安防领域，项目建设深圳市域级三维多模态感知平台，并形成电动车监管平台和充电桩管理平台等应用。相关平台已接入物联网设备超过300万台，在1年内发现违规行为2.3万次，有效预警火灾约150次，挽回经济损失超过1亿元。通过对分布式设备状态、空间信息和风险事件进行统一接入与综合分析，项目技术提高了城市安全风险的主动发现能力，为街道和相关管理部门开展电动车监管、充电安全管理及火灾隐患处置提供了技术支撑。

超过300万台物联网设备形成了高并发、持续更新的城市级感知网络。项目平台需要同时处理设备状态、位置关系、告警信息和历史记录，并从大量常规数据中识别少量但具有较高风险的异常事件。多源协同融合和可信采信方法在这一场景中体现为对不同终端信息的统一接入、关联分析与风险筛选，使管理人员能够把注意力集中到需要核查和处置的事件上。

违规行为发现、火灾预警和经济损失挽回分别反映了平台从过程监管到风险阻断的不同层次价值。发现2.3万次违规行为说明系统能够形成常态化监管能力；约150次有效火灾预警说明感知结果能够进一步服务安全事件识别；挽回经济损失超过1亿元则体现了提前发现与及时处置产生的风险避免效益。该类效益虽然不同于产品销售形成的直接收入，却是城市公共安全应用经济价值的重要组成部分。

\subsubsection{智慧消防与应急管理}

在智慧应急领域，项目成果已规模化应用于广东省多地消防救援支队和基层消防单位，形成深圳市消防救援支队智慧平台、石岩街道水田新村智慧消防系统等典型应用。平台累计接入设备终端超过5万台，上线半年内发现约230起报警，故障切换时间小于1秒，累计接入感知数据24亿条。相关应用通过设备感知、报警汇聚、异常识别和快速切换，增强了消防安全状态的连续监测与应急响应能力，为市级消防平台与基层消防单位之间的信息协同提供了支撑。

智慧消防应用同时覆盖市级指挥平台和基层社区系统。市级平台承担大范围消防信息汇聚、综合展示和协同管理，基层系统则直接连接具体设备与现场风险，两类系统通过统一的数据接入和报警机制形成上下联动。超过5万台终端和24亿条感知数据说明平台已经经历较大规模、较长时间的数据运行检验，而半年内发现约230起报警则反映了平台从设备接入走向业务应用的实际成效。

对于消防应急系统而言，设备或链路故障可能影响告警信息的及时传递，因此故障切换时间小于1秒具有直接的业务意义。该指标与报警发现数量共同说明，项目应用不仅关注感知准确性，也将平台连续运行和异常情况下的快速恢复纳入系统建设。相关消防救援支队和基层消防单位出具的应用材料进一步证明，成果已经进入真实消防安全管理流程，而非停留在演示性部署阶段。

\subsection{规模化推广基础与可复制性}

项目成果具备规模化推广的首要基础，是已经形成覆盖不同数据规模和业务层级的应用验证。智能制造以单条产线和检测工位为基本单元，可以通过复制算法模块和设备接口扩展到更多生产线；智慧交通以道路、桥梁、隧道、边坡和枢纽为管理对象，可以在统一平台上持续增加设施类型和覆盖区域；智能安防与智慧应急则以物联网终端和业务单位为扩展单元，可以在保留核心平台的基础上增加设备、街道或消防单位。不同场景均已呈现出由单点接入向多节点、由局部业务向综合平台扩展的特征。

第二项基础是项目技术具有分层组合的产品结构。感知算法可以作为独立模块嵌入既有设备，多源融合和数据治理能力可以作为平台中间层连接不同终端，高层认知和任务协同能力可以根据业务需要配置在管理与决策端。这种结构允许用户从最迫切的检测或监测需求开始建设，再逐步扩展到跨设备融合和综合业务协同，既降低初始部署门槛，也为后续系统升级保留空间。

第三项基础是项目已经形成跨行业用户和应用证明。交通运输管理部门、消防救援支队、街道基层单位以及制造企业从不同角度验证了成果的可用性：交通用户验证大范围基础设施监测能力，制造企业验证精细质检和经营收益，安防与消防用户验证海量终端接入、异常预警和连续运行能力。这些应用结果为项目向同类型地区和企业推广提供了事实依据，也有助于将用户需求进一步转化为产品规范、实施流程和运维经验。

从后续推广方式看，可以形成以行业标杆应用带动同类场景复制、以平台能力复用带动区域扩展、以算法模块升级带动存量系统增值的多层次模式。对制造企业，可围绕不同缺陷类型和工艺环节增加检测模块；对交通管理部门，可在既有平台上扩展设施覆盖和风险分析功能；对安防与应急用户，可进一步增加终端类型、预警规则和协同处置流程。上述推广方式不需要改变项目技术主线，能够延续现有成果的工程积累和应用基础。

\subsection{综合经济与社会效益}

项目成果的经济效益不仅体现为产品销售、平台建设和技术服务带来的直接收益，也体现为智能检测降低质量损失、设施监测减少运维风险、灾害预警避免财产损失等间接价值。智能制造应用形成了超过66亿元的直接经济效益，智能安防应用挽回经济损失超过1亿元；与此同时，交通监测和智慧消防等平台通过提高风险发现的及时性与处置效率，为用户降低潜在事故和设备故障带来的运营成本。

从直接经济效益看，项目通过智能制造生产线建设、检测系统部署、城市感知平台建设和相关技术服务，将科研成果转化为应用单位的销售收入与新增效益。制造场景中的效益可以通过产线数量、产品良率、销售增量和直接经济效益进行量化，交通平台建设也形成了明确的销售收入和直接经济效益。这部分价值说明项目成果已经形成可交易、可交付和可持续运行的产品或服务，而不是一次性的科研演示。

从间接经济效益看，交通设施的在线监测有助于减少重复巡查和分散管理成本，智能质检有助于降低漏检、返工和质量波动造成的损失，安防与消防预警有助于在风险扩大前启动处置。现有材料已量化的挽回经济损失超过1亿元，是间接效益的代表性结果；其他安全运行、运维优化和处置提效价值由于缺少统一货币化口径，本报告不再作额外金额推算。通过区分已经统计的直接收益与尚未货币化的风险避免价值，可以避免重复计算，同时更加完整地呈现成果的实际贡献。

在社会效益方面，项目成果服务于交通基础设施安全、制造业质量提升、城市基层治理和消防应急保障等公共需求。大规模交通设施监测提升了道路、桥梁、隧道和边坡的数字化管养能力；数百万级物联网设备接入增强了城市安防覆盖范围；多地消防救援支队和基层消防单位的应用推动了消防信息的实时汇聚与协同处置。项目由此形成了“技术研发---平台建设---行业应用---治理增效”的转化路径，在创造直接经济效益的同时，为城市安全运行、产业数字化升级和公共服务能力提升提供了支撑。

具体而言，智慧交通应用把分散设施纳入城市级运行监测，提升了基础设施状态的可见性和管理协同性；智能制造应用推动精细质量检测由人工经验向自动化、数字化和智能化转变，服务制造业质量升级；智能安防应用通过常态化发现违规行为和预警火灾，增强基层治理的主动性；智慧应急应用连接市级消防平台与基层消防单位，提高报警信息汇聚和系统故障恢复能力。这些社会效益共同指向更安全、更高效和更具韧性的城市运行体系。

总体来看，项目已经实现从核心技术突破到规模化产业应用的跨越：在产业端形成1379.52亿元销售收入和72.29亿元直接经济效益，在应用端覆盖43条生产线、大规模交通基础设施、300万台以上安防物联网设备和5万台以上消防设备终端，在治理端形成交通设施监测、违规行为发现、火灾预警与消防协同等实际能力。上述结果从经营收益、应用规模和公共价值三个维度证明了项目成果的转化成效，并为进一步扩大行业覆盖和区域推广提供了基础。
\clearpage
